\begin{figure}[h]
  \centering

  \begin{tikzpicture}[>=Stealth, scale=1.3]

    % --- 1. TOTAL SPACE TM (Top Left) ---
    \begin{scope}[shift={(-3.5,0)}]
      % Path definitions for the 3D "out-in-out" effect
      \def\Mcurve{(-2,0.2) to[out=20, in=160] (0,-0.1) to[out=-20, in=200] (2,0.2)}
      \def\TMtop{(-2,1) to[out=20, in=160] (0,0.7) to[out=-20, in=200] (2,1)}
      \def\TMbottom{(-2,-0.6) to[out=20, in=160] (0,-0.9) to[out=-20, in=200] (2,-0.6)}
      \def\TMpath{\TMtop -- (2,-0.6) -- \TMbottom -- (-2,-0.6) -- cycle}

      % Shading the local patch pi^-1(U)
      \begin{scope}
        \clip \TMpath;
        \fill[gray!30] (-1,-1.5) rectangle (1,1.5);

        % Vertical Fibers (Tangent Spaces T_pM)
        \foreach \x in {-1.9, -1.7, ..., 1.9} {
          \draw[gray!40, thin] (\x, -1.5) -- (\x, 1.5);
        }
      \end{scope}

      % Dashed boundaries for the patch
      \begin{scope}
        \clip \TMpath;
        \draw[thick, dashed] (-1,-1.5) -- (-1,1.5);
        \draw[thick, dashed] (1,-1.5) -- (1,1.5);
      \end{scope}

      % The Zero Section (s: M -> TM)
      \draw (-2, 0.2) to[out=20, in=160] (0, -0.1) to[out=-20, in=200] (2, 0.2);

      \draw[thick] \Mcurve;

      % Overlay thick line ONLY inside the patch area using clipping
      \begin{scope}
        \clip (-1,-1) rectangle (1,1);
        \draw[ultra thick] \Mcurve;
      \end{scope}

      \node[above] at (0, 1.1) {\large $TM$};
      \node[below] at (0.8, -1) {\large $\pi^{-1}(U)$};
    \end{scope}

    % --- 2. TARGET SPACE R^2n (Top Right) ---
    \begin{scope}[shift={(1.5,0)}]
      \filldraw[fill=gray!30, draw=white] (-1,-0.8) rectangle (1,0.8);

      \foreach \x in {-0.8, -0.6, ..., 0.8} {
        \draw[black, thin] (\x, -0.8) -- (\x, 0.8);
      }
      \draw[black, dashed] (-1, -0.8) -- (-1, 0.8);
      \draw[black, dashed] (1, -0.8) -- (1, 0.8);

      \draw[<->, ultra thick] (-1.8, 0) -- (1.8, 0);

      \node[above] at (0, 0.8) {\large $\mathbb{R}^{2n}$};
      \node[below] at (0, -0.8) {\large $\varphi(U) \times \mathbb{R}^n$};
    \end{scope}

    % --- 3. BASE MANIFOLD M (Bottom Left) ---
    \begin{scope}[shift={(-3.5,-3.5)}]
      % Define the curve base path to match TM bottom exactly
      \def\Mcurve{(-2,0.2) to[out=20, in=160] (0,-0.1) to[out=-20, in=200] (2,0.2)}

      % Draw regular thickness curve outside/inside
      \draw[thick] \Mcurve;

      % Overlay thick line ONLY inside the patch area using clipping
      \begin{scope}
        \clip (-1,-1) rectangle (1,1);
        \draw[ultra thick] \Mcurve;
      \end{scope}

      % Endpoint parentheses
      \node at (-0.98,0.2) {\huge (};
      \node at (0.98,-0.2) {\huge )};

      \node at (2.3, 0.3) {\large $M$};
      \node at (0, -0.4) {\large $U$};
    \end{scope}

    % --- 4. COORDINATE IMAGE phi(U) (Bottom Right) ---
    \begin{scope}[shift={(1.5,-3.5)}]
      \draw[<->, thick] (-1.8, 0.05) -- (1.8, 0.05);
      \draw[ultra thick] (-1, 0.05) -- (1, 0.05) node[midway, below] {\large $\phi(U)$};
      \draw[ultra thick] (-0.98, 0.05) node {\huge (} -- (0.98, 0.05) node {\huge )};
    \end{scope}

    % --- 5. MAPPING ARROWS ---
    \draw[->, thick, orange!80!black] (-1, 0.4) to[bend left=25] node[midway, above] {\huge $\widetilde{\varphi}$} (0.5, 0.4);
    \draw[->, thick, blue!70!black] (-2.8, -3.5) to[bend left=25] node[midway, above] {\huge $\varphi$} (1.3, -3.1);

    \draw[->, thick] (-3.5, -1.2) -- (-3.5, -2.8) node[midway, left] {\huge $\pi$};
    \draw[->, thick] (2, -1.2) -- (2, -2.8);

  \end{tikzpicture}

  \caption{Coordinates for the Tangent bundle}
  \label{fig:coordinates_for_the_tangent_bundle}
\end{figure}
